\paragraph{\trArnoldStratification。}\label{sub:arnold-stratification}
积 $Gr(V)^2=Gr(V)\times Gr(V)$ 带有一个按相对位置给出的%
\textbf{\trSchubert{}型\trStratum}（Schubert-type stratification）。每个\trStratum{}
$\Sigma^d=\{(\Lambda_0,\Lambda_1):\dim(\Lambda_0\cap\Lambda_1)=d\}$ 都带着一个由这些交构
成的、秩为 $d$ 的\trVectorBundle{} $\sigma^d$，且配有\trCanonical{}\trIsomorphism{}
\begin{equation}\label{eq:normal-bundle-isomorphism}
(\mathrm{sym}^2\sigma^d)^\vee\cong\nu_{\Sigma^d}.
\end{equation}
可以把\cref{sub:path-spaces-spectral-flow} 中的\trPathSpace{}构造应用于 $Gr(V)^2$，取
$\Sigma=\Sigma^1$、$\sigma=\sigma^1$。其结果是一个\trLocallyConstant{}的整值函数
$I_Z$，即\textbf{路径的\trMaslovIndex}（Maslov index for paths），以及一个%
\trRealLineBundle{} $\delta_Z$；二者都住在由路径
$(\lambda_0,\lambda_1)\colon[0;1]\to Gr(V)^2$——满足 $\lambda_0(s)\pitchfork\lambda_1(s)$
（$s=0,1$）——构成的\trPathSpace{} $\scr{P}^2\,Gr(V)$ 上。一如先前考虑过的%
\trQuadraticForm{}情形——后者实际上可以充当局部模型——这里有一个易于克服的复杂
之处，涉及 $\Sigma^2$。
值得把这一点说得稍微显式些。取一条路径 $(\lambda_0,\lambda_1)$ 与某个 $s\in[0;1]$。
选取\trLinear{}映射的族 $\phi_{k,r,s}\colon\lambda_k(s)\to\lambda_k(r)$（$k=0,1$，
$|r-s|$ 充分小），使得 $\phi_{k,s,s}=\mathrm{Id}$。\trCrossingForm{}是
$\lambda_0(s)\cap\lambda_1(s)$ 上由下式定义的\trQuadraticForm：
\begin{equation*}
q_{\lambda_0,\lambda_1}(s)\colon v\longmapsto
-\Bigl(\frac{d}{dr}\Bigr)_{\!r=s}\omega_V\bigl(\phi_{0,r,s}(v),\phi_{1,r,s}(v)\bigr).
\end{equation*}
它不依赖于诸 $\phi_{k,r,s}$。用\eqref{eq:normal-bundle-isomorphism} 的话说：若对某个
$d>0$ 有 $(\lambda_0(s),\lambda_1(s))\in\Sigma^d$，则\trCrossingForm{}是对应于
$(\lambda_0'(s),\lambda_1'(s))$ 的法向分量的\trQuadraticForm。对路径的一个\trGeneric{}
选取，诸空间 $\lambda_0(s)\cap\lambda_1(s)$ 至多是一维的，而一维空间上的诸%
\trCrossingForm{}将是非零的。于是，应用一般机制\eqref{eq:delta-alpha} 便可得到
$I_Z(\lambda_0,\lambda_1)$ 与 $(\delta_Z)_{\lambda_0,\lambda_1}$。
与\cref{sub:path-spaces-spectral-flow} 相平行，我们现在来解释这些\trInvariant{}在%
\trIndex{}理论上的含义。照该处的讨论选取一个函数 $\psi$。对每个
$(\lambda_0,\lambda_1)\in\scr{P}^2\,Gr(V)$，关联一个算子 $D_{Z,\lambda_0,\lambda_1}$：
它是无限\trStrip{}上平凡\trVectorBundle{} $E=Z\times V$ 的 $\bar\partial$-算子，其%
\trBoundaryCondition{}由\trLagrange\trSubbundle{} $F_{\lambda_0,\lambda_1}\subset E|\partial Z$
给出：$(F_{\lambda_0,\lambda_1})_{(s,k)}=\lambda_k(\psi(s))$。由于端点处的%
\trTransversality{}条件，这是一个\trNondegenerate{}的\trCauchyRiemann{}算子，从而是%
\trFredholm{}算子。于是
\begin{lemma}\label{thm:path-space-det-line}
$I_Z(\lambda_0,\lambda_1)=\mathrm{index}(D_{Z,\lambda_0,\lambda_1})$；且
$(\delta_Z)_{\lambda_0,\lambda_1}$ 与 $\det(D_{Z,\lambda_0,\lambda_1})$
\trCanonical{}\trIsomorphic。
\end{lemma}
对此的策略与\cref{thm:spectral-flow-index} 之证明的第二版相同，故我们不再重复整个
论证。代替给出\eqref{eq:cut-decomposition} 的初等常\trDifferential{}方程计算，这里直接
使用\trGluing{}公式\eqref{eq:delbar-gluing}。剩下的任务是找到足够多的例子，使%
\trIndex{}与\trDeterminantLine{}能被显式地写出。下面是一种做法（还可以有别的做法，其中
一些在技术上略微简单些）。设 $V=\RR^{2n}=\CC^n$ 带标准\trSymplecticStructure{}与
标准\trComplexStructure{} $I_V=i$；$\lambda_1(s)=\RR^n\times\{0\}$ 对一切 $s$；而
$\lambda_0(s)$ 是一个图 $\{x_2=\alpha(s)x_1\}$，其中 $\alpha(s)$ 是对称 $n\times n$ 矩阵
的一个族。\trEquivalence{}地说，这意味着 $\lambda_0(s)=\exp(iA(s))\RR^n$，其中
$A(s)=\begin{psmallmatrix}\alpha(s) & 0\\ 0 & 0\end{psmallmatrix}$。如果随后通过施加辛%
\trLinear{}映射的族 $\Phi(s,t)=\exp((1-t)\,iA(\psi(s)))$（$(s,t)\in Z$）来改变 $E$ 的%
\trTrivialization，该算子就变成
\begin{equation}\label{eq:deformed-cr-operator}
D_{Z,\lambda_0,\lambda_1}X=(\partial_s-J(s,t)iA(\psi(s)))X+J(s,t)\partial_t X
+\Phi(s,t)^{-1}(\partial_s\Phi(s,t))X,
\end{equation}
其中 $J(s,t)=\Phi(s,t)^{-1}i\Phi(s,t)$ 是 $\RR^{2n}$ 上一族\trCompatible{}%
\trComplexStructure。可以把 $J(s,t)$ 经连续\trDeformation{}带回常值族 $i$；这就给出
\eqref{eq:deformed-cr-operator} 的一个经\trFredholmOperator{}的\trHomotopy，因而%
\trIndex{}与（在\trCanonical{}\trIsomorphism{}意义下的）\trDeterminantLine{}保持不变。
类似地，\eqref{eq:deformed-cr-operator} 中的最后一项是\trCompactlySupported{}且零阶的，
从而是紧的，可以经\trDeformation{}把它去掉。所得的算子是
\begin{equation*}
\tilde{D}_Z X=(\partial_s+A(\psi(s)))X+i\partial_t X
\end{equation*}
带常值\trBoundaryCondition{} $X(s,0)$、$X(s,1)\in\RR^n\times\{0\}$。一个%
\trEllipticRegularity{}论证表明：\trKernel{}的每个元素都可延拓为 $\RR^2$ 上一个满足%
\trPeriodic{}性 $X(s,t+2)=X(s,t)$、$X(s,-t)=\overline{X}(s,t)$ 的\trSmooth{}函数；对%
\trCokernel{}同理。把\trDomain{}与值域限制到这些子空间之后，便可沿 $t$ 方向施加%
\trFourier{}变换，这将 $\tilde{D}_Z$ \trSplitting{}成如下形状的一维%
\trDifferential{}算子的\trDirectSum：
\begin{equation}\label{eq:fourier-summands}
d/ds+\alpha(\psi(s))
\qquad\text{以及}\qquad
d/ds+\begin{pmatrix}\alpha(\psi(s)) & \pi l\cdot\mathrm{Id}\\ \pi l\cdot\mathrm{Id} & 0\end{pmatrix}
\quad\text{其中 } l=1,2,\ldots.
\end{equation}
这类算子的\trIndex{}理论我们先前已在\cref{sub:path-spaces-spectral-flow} 中研究过。对
\eqref{eq:fourier-summands} 而言，结论是：唯一的非平凡贡献来自第一个被加项，更具体
地说，来自那些使 $\alpha(s)$ 的\trNullSpace{}非空的 $s$。一个简单的计算表明，这些贡献
恰好是 $(\lambda_0,\lambda_1)$ 的诸\trCrossingForm。
\begin{remark}\label{thm:arnold-stratification-references}
\trLagrangianGrassmannian{}的这一\trStratum{}是在 [2] 中引入的。%
\cref{thm:path-space-det-line} 的第一部分是\trFloer{}的\trIndex{}定理 [45]。总的说来，
我们的叙述紧密遵循 [112]、[113]。上面的最后一个论证可以看作从\trFloer{}理论到%
\trMorse{}理论的化归 [47] 的一个\trLinear{}版本。
\end{remark}
